\documentclass[12pt]{article}
\usepackage[utf8]{inputenc}
\usepackage[T1]{fontenc}
\usepackage{hyperref}
\title{Diuxi-Tilantongo Mixtec Language Resources for K-12 Curriculum Development}
\author{Compiled by Claude Sonnet 5.5 (claude-sonnet-5-5) for Andr\'as Kornai}
\date{}
\begin{document}
\maketitle

\section{Introduction}
Diuxi-Tilantongo Mixtec (ISO 639-3 xtd; Glottocode diux1235; also Central Nochixtl\'an Mixtec, Tnu'un dau) is an Otomanguean language of Oaxaca, Mexico. Wikipedia notes that numbers are declining through emigration to the United States, and Wikitongues rates it 6b (Threatened). It has more material than the other three languages in this batch, though most of it is SIL-era or religious.

\section{Primary Educational Resources}
\begin{itemize}
\item \emph{Cartilla mixteca}, a 109-page primer for the Mixteco of San Juan Diuxi and Santiago Tilantongo, published by the Summer Institute of Linguistics and held in the El Colegio de M\'exico repository: \url{https://repositorio.colmex.mx/concern/books/cz30px32x?locale=en}
\item Wikipedia also mentions a Spanish-language transition primer, \emph{?Na ka'u-ro tnu'u \~nudau!} (2006). Its own link was not checked, and the Wikipedia page is: \url{https://en.wikipedia.org/wiki/Tilantongo_Mixtec}
\end{itemize}

\section{Linguistic and Reference Materials}
\begin{itemize}
\item Wikipedia article listing the SIL-era literature, including Oram's \emph{A sketch of Mixtec grammar} (1970), Pike and Oram on stress and tone (1976), and Kuiper and Oram's \emph{A Syntactic Sketch of Diuxi-Tilantongo Mixtec}: \url{https://en.wikipedia.org/wiki/Tilantongo_Mixtec}
\item Wikitongues language page: \url{https://wikitongues.org/languages/xtd/}
\end{itemize}

\section{Religious and Community Resources}
\begin{itemize}
\item Scripture Earth lists a Bible in this language, with a New Testament PDF and an Android app: \url{https://www.scriptureearth.org/00eng.php?iso=xtd}
\end{itemize}

\section{Digital and Audio Resources}
\begin{itemize}
\item Global Recordings Network audio Bible lessons and short evangelistic programmes, with a downloadable audio zip: \url{https://globalrecordings.net/en/language/xtd}
\item Second Global Recordings Network listing (Diuxi-Tilantongo Mixtec language): \url{https://globalrecordings.net/en/language/6126}
\end{itemize}

\section{Curriculum Development Context}
A SIL primer, grammar sketches and a New Testament provide a base, so an orthography exists, but none of this is a school curriculum. The materials are decades old and largely religious. Any programme would need community review and new content for non-religious subjects.
\end{document}